\section*{Arithmetic formalization of the bilateral balance}
\addcontentsline{toc}{section}{Arithmetic formalization of the bilateral balance}
\label{libro:formalizacion-lean}

The nonadic specialization of the bilateral law admits an independent
arithmetic formalization. The point of application comes after the
construction of its coefficients and transports: the base-ten partition fixes
\(a=8\), and the two registers are \(8x-y\) and \(9x+y\). On this domain,
Lean~4.21.0 verifies seven theorems using its \texttt{Std} library. The formal
sources and execution procedure are provided as a supplement.
This addendum states their mathematical content and relates it to
the operator proofs in the article.

\subsection*{Polynomial identity and conservation of the norm}

For all \(x,y\in\mathbb Z\), we have
\[
 9(8x-y)^2+8(9x+y)^2=17(72x^2+y^2).
\]
Expanding the two squares cancels the cross terms. When
\(y^2=x^2\), the same equality gives
\[
 9(8x-y)^2+8(9x+y)^2=17\cdot73\,x^2.
\]
These are the first two formalized theorems. The first expresses a
balance for arbitrary integer inputs; the second applies the
explicit norm-conservation hypothesis. The integer \(73=72+1\) occupies
exactly the place it has in \cref{lib04:teo-balance}.

\subsection*{Arbitrary finite dimension}

Let \(n\in\mathbb N\) and \(x,y:\mathbb N\to\mathbb Z\). Define
\[
 S_n(x)=\sum_{j<n}x_j^2,
 \qquad
 E_n(x,y)=\sum_{j<n}\bigl(9(8x_j-y_j)^2+8(9x_j+y_j)^2\bigr).
\]
The next two theorems establish
\[
 E_n(x,y)=17\bigl(72S_n(x)+S_n(y)\bigr)
\]
and, under the hypothesis \(S_n(y)=S_n(x)\),
\[
 E_n(x,y)=17\cdot73\,S_n(x).
\]
The formal proof proceeds by induction on \(n\). The initial case is the
empty sum; the induction step adds one coordinate and applies the polynomial
identity already proved. The universal quantifier over \(n\) includes
all finite dimensions. The hypothesis uses the total sum of
squares, so it allows redistributions among coordinates.

\subsection*{Recovery of the inputs}

Write \(q_-=8x-y\) and \(q_+=9x+y\). The fifth and sixth theorems
are the exact equalities
\[
 q_-+q_+=17x,\qquad 8q_+-9q_-=17y.
\]
On the image of the map \((x,y)\mapsto(q_-,q_+)\), both left-hand
sides are multiples of \(17\). Division recovers the two
inputs. The seventh theorem states their injectivity consequence:
\[
 \left.
 \begin{aligned}
 8x-y&=8x'-y',\\
 9x+y&=9x'+y'
 \end{aligned}
 \right\}
 \quad\Longrightarrow\quad x=x'\ \text{and}\ y=y'.
\]
Cancellation in \(\mathbb Z\) justifies this conclusion. The formal
content coincides with recovery of the pair of registers before
they are normalized as channels of a metric realization.

\subsection*{Relation between the proofs}

The formalization uses arbitrary integer variables and mathematical
induction. It therefore constitutes a proof of the preceding statements,
distinct from checking a finite set of examples. The system
accepts the proof terms without added axiomatic declarations or
deferred proofs. The supplement also records the logical
dependencies reported by the tool for each declaration.

The operator law in the body of the article has a broader domain:
isometries between Hilbert spaces, an archive of arbitrary depth and
conjugate readers. Its proofs are found in
\cref{lib04:teo-balance,lib06:simetria,lib06:profundidad}. The Lean proof
included here certifies the stated arithmetic specialization, with
its quantifiers and starting structure. The provenance of the
coefficients, generation of the register and exceptional realizations
retain the earlier proofs on which they depend.
